\section{Robustness of cover removal and the timing-matched countermeasure}\label{supp:c-ext}
This section tests three objections to the removal attack of Section~\ref{supp:b1-timing}: that it needs cover labels, that it would not survive network timing noise, and that a defender could simply delay its decoys. All three analyses reuse the collected 3,600 episodes and their 120/60/120 split. They involve no new LLM requests or captures.

\paragraph{Freeze and exposure}
The plan, code and input hashes were frozen and committed before any of their test quantities was computed. Before the freeze, the removal results of Section~\ref{supp:b1-timing} were known. Validation gap extremes (cover at most 56\,ms, honest at least 418\,ms) and that section's validation candidate table had also been viewed. A reproduction gate refitted the Section~\ref{supp:b1-timing} pipeline first. It required, and obtained, exact equality with the original size-view LR/RF test predictions and with all twelve of its strip counts. An independent script recomputes the label-free detector counts from raw receipts, every McNemar and Holm value, and the summary criteria.

\paragraph{Label-free removal}
The attacker fits a two-component Gaussian mixture to the log gaps of unlabeled training calls. The threshold is the point of equal posterior responsibility: 0.189\,s for Qwen and 0.142\,s for Llama. Table~\ref{tab:c2} shows perfect test detection in every cell and strip recognition equal to or above the supervised detector.

\paragraph{Timestamp jitter}
Independent half-normal delays are added to every call start and end in all splits, and gaps and durations are recomputed (five draws per $\sigma$). A detector frozen on clean data degrades as $\sigma$ grows; a detector refitted on jittered training data does not, until $\sigma=200$\,ms reduces Qwen's recall to about .87 (Table~\ref{tab:c1}). Jitter is independent per event. Correlated queueing delay, which shifts neighbouring events together, would mostly preserve gaps. The analysis does not model packetization or loss.

\paragraph{Timing-matched dispatch}
The counterfactual gateway delays each decoy by a gap drawn from the defender's own no-cover training gaps (\emph{gap-matched}), and in the \emph{gap+duration-matched} variant also extends its duration to a draw from no-cover durations at the same destination. Each added delay shifts all later timestamps of the episode. Hidden cover never reaches the planner and model inference does not depend on wall-clock time, so the planner's decisions are unchanged. Visible cover changes only timing, not content. The attacker refits both the detector and the direct timing-plus-size classifier on transformed data. Table~\ref{tab:c3} reports every cell. Residual separability remains (AUROC .74--.88), but the detector's precision of .35--.58 makes stripping counter-productive. The pre-registered summary rule asked whether gap+duration-matched strip recognition falls within ten points of the size view in every cell. It does not: in 7 of 10 cells it falls more than ten points \emph{below}. The best attack reported in this section (byte learner, timing-plus-size and strip) is within three tasks of the byte learner in every cell. That comparison omits position rules and attacks that adapt to the matched cover; Section~\ref{supp:d-ext} adds both and revises the conclusion.

\begin{table*}[t]\centering\small
\caption{Label-free removal. The threshold is the equal-responsibility point of a two-component Gaussian mixture fitted to unlabeled training log-gaps. TP/FP/FN count test calls (first ticket call excluded). Difference is Strip minus Byte in percentage points with a conditional 95\% interval; one ten-test Holm family. Fixed 0.1\,s and 0.2\,s thresholds give identical counts.}\label{tab:c2}
\setlength{\tabcolsep}{4pt}
\begin{tabularx}{\textwidth}{@{}llrrrrYr@{}}\toprule
Model & Policy & Thr.\ (s) & TP/FP/FN & Byte & Strip & Strip$-$Byte [95\% CI] & Holm $p$\\\midrule
Qwen3:8B & Proxy discard & 0.189 & 182/0/0 & 62 & 115 & +44.2 [35.0, 53.3] & 3.1e-14 \\
Qwen3:8B & Read & 0.189 & 182/0/0 & 62 & 114 & +43.3 [34.2, 52.5] & 4.9e-14 \\
Qwen3:8B & Isolated, hidden & 0.189 & 182/0/0 & 61 & 114 & +44.2 [35.0, 53.3] & 3.1e-14 \\
Qwen3:8B & Live write & 0.189 & 182/0/0 & 62 & 112 & +41.7 [31.7, 51.7] & 1.8e-11 \\
Qwen3:8B & Isolated, visible & 0.189 & 182/0/0 & 63 & 111 & +40.0 [30.8, 49.2] & 1.8e-11 \\
Llama3.1:8B & Proxy discard & 0.142 & 182/0/0 & 79 & 116 & +30.8 [23.3, 39.2] & 5.8e-11 \\
Llama3.1:8B & Read & 0.142 & 182/0/0 & 72 & 112 & +33.3 [25.0, 41.7] & 1.3e-11 \\
Llama3.1:8B & Isolated, hidden & 0.142 & 182/0/0 & 82 & 112 & +25.0 [17.5, 33.3] & 4.6e-08 \\
Llama3.1:8B & Live write & 0.142 & 182/0/0 & 80 & 112 & +26.7 [18.3, 35.8] & 1.3e-07 \\
Llama3.1:8B & Isolated, visible & 0.142 & 182/0/0 & 81 & 107 & +21.7 [13.3, 30.0] & 6.2e-06 \\
\bottomrule\end{tabularx}
\end{table*}

\begin{table*}[t]\centering\small
\caption{Jitter: strip recognition (correct/120, mean of five jitter draws) under independent half-normal timestamp jitter. Each cell gives refitted / frozen-clean detector. Byte is the size-view selected attack. The pre-registered robustness criterion (refitted mean at least 24 above Byte in every cell) holds at every $\sigma$.}\label{tab:c1}
\setlength{\tabcolsep}{4pt}
\begin{tabularx}{\textwidth}{@{}llYrrrr@{}}\toprule
Model & Policy & Byte & $\sigma$=10\,ms & 50\,ms & 100\,ms & 200\,ms\\\midrule
Qwen3:8B & Proxy discard & 62 & 114.8 / 113.6 & 114.0 / 99.8 & 113.2 / 85.6 & 103.8 / 79.4 \\
Qwen3:8B & Read & 62 & 113.8 / 112.4 & 113.2 / 94.6 & 113.2 / 83.4 & 103.4 / 77.4 \\
Qwen3:8B & Isolated, hidden & 61 & 112.4 / 111.0 & 114.0 / 93.4 & 112.2 / 86.6 & 105.6 / 80.0 \\
Qwen3:8B & Live write & 62 & 110.8 / 107.6 & 111.6 / 88.6 & 111.0 / 85.0 & 102.0 / 77.8 \\
Qwen3:8B & Isolated, visible & 63 & 110.2 / 108.6 & 110.6 / 92.2 & 109.8 / 79.2 & 98.4 / 76.8 \\
Llama3.1:8B & Proxy discard & 79 & 115.6 / 115.4 & 115.8 / 105.2 & 116.0 / 95.0 & 114.8 / 90.6 \\
Llama3.1:8B & Read & 72 & 112.0 / 112.0 & 111.4 / 102.4 & 112.0 / 91.4 & 109.8 / 85.6 \\
Llama3.1:8B & Isolated, hidden & 82 & 112.0 / 112.0 & 111.8 / 98.6 & 112.0 / 93.2 & 110.0 / 85.6 \\
Llama3.1:8B & Live write & 80 & 112.0 / 112.0 & 112.0 / 98.6 & 111.8 / 96.6 & 110.4 / 86.8 \\
Llama3.1:8B & Isolated, visible & 81 & 107.0 / 107.0 & 107.0 / 97.2 & 107.0 / 91.4 & 105.4 / 83.2 \\
\bottomrule\end{tabularx}
\end{table*}

\begin{table*}[t]\centering\small
\caption{Timing-matched dispatch (receipt-level counterfactual). GM draws each decoy gap, GDM its gap and duration, from no-cover training history; the attacker is fully refitted. Detector P/R/AUROC on test calls. Unm.\ is the unmatched supervised strip score (Section~\ref{supp:b1-timing}), Strip the matched strip score, T+S the refitted direct timing-plus-size attack, Byte the size view. Latency is mean added seconds per covered episode (percentage of no-cover wall time). Holm columns: GDM Strip vs.\ Unm.\ (primary family) and GDM Strip vs.\ Byte (secondary family).}\label{tab:c3}
\setlength{\tabcolsep}{4pt}
\begin{tabularx}{\textwidth}{@{}llrYrrrrrrr@{}}\toprule
Model & Policy & Var. & P/R/AUC & Unm. & Strip & T+S & Byte & +Latency s & Holm (Unm.) & Holm (Byte)\\\midrule
Qwen3:8B & Proxy discard & GM & 0.57/0.66/0.80 & 114 & 30 & 57 & 62 & 1.07 (12.3\%) & -- & -- \\
Qwen3:8B & Read & GM & 0.58/0.67/0.80 & 114 & 30 & 52 & 62 & 1.05 (12.0\%) & -- & -- \\
Qwen3:8B & Isolated, hidden & GM & 0.57/0.65/0.80 & 112 & 31 & 47 & 61 & 1.07 (12.3\%) & -- & -- \\
Qwen3:8B & Live write & GM & 0.57/0.66/0.79 & 111 & 30 & 61 & 62 & 1.04 (11.9\%) & -- & -- \\
Qwen3:8B & Isolated, visible & GM & 0.57/0.66/0.79 & 111 & 32 & 55 & 63 & 1.04 (11.9\%) & -- & -- \\
Qwen3:8B & Proxy discard & GDM & 0.54/0.79/0.86 & 114 & 43 & 57 & 62 & 1.08 (12.4\%) & 1.5e-18 & 0.174 \\
Qwen3:8B & Read & GDM & 0.55/0.79/0.87 & 114 & 50 & 49 & 62 & 1.09 (12.4\%) & 1.4e-16 & 0.495 \\
Qwen3:8B & Isolated, hidden & GDM & 0.57/0.77/0.88 & 112 & 57 & 41 & 61 & 1.12 (12.8\%) & 9.9e-14 & 0.733 \\
Qwen3:8B & Live write & GDM & 0.56/0.81/0.87 & 111 & 54 & 54 & 62 & 1.11 (12.7\%) & 1.5e-15 & 0.733 \\
Qwen3:8B & Isolated, visible & GDM & 0.53/0.76/0.86 & 111 & 49 & 66 & 63 & 1.08 (12.3\%) & 4.8e-16 & 0.462 \\
Llama3.1:8B & Proxy discard & GM & 0.38/0.68/0.75 & 115 & 48 & 72 & 79 & 1.20 (14.7\%) & -- & -- \\
Llama3.1:8B & Read & GM & 0.38/0.69/0.77 & 112 & 44 & 69 & 72 & 1.21 (14.8\%) & -- & -- \\
Llama3.1:8B & Isolated, hidden & GM & 0.39/0.73/0.76 & 112 & 52 & 84 & 82 & 1.22 (14.9\%) & -- & -- \\
Llama3.1:8B & Live write & GM & 0.39/0.71/0.77 & 112 & 55 & 66 & 80 & 1.21 (14.8\%) & -- & -- \\
Llama3.1:8B & Isolated, visible & GM & 0.35/0.61/0.74 & 107 & 48 & 79 & 81 & 1.23 (15.1\%) & -- & -- \\
Llama3.1:8B & Proxy discard & GDM & 0.46/0.65/0.82 & 115 & 59 & 79 & 79 & 1.33 (16.4\%) & 1.6e-15 & 0.070 \\
Llama3.1:8B & Read & GDM & 0.49/0.70/0.84 & 112 & 59 & 74 & 72 & 1.29 (15.8\%) & 1.5e-15 & 0.495 \\
Llama3.1:8B & Isolated, hidden & GDM & 0.46/0.65/0.83 & 112 & 55 & 81 & 82 & 1.28 (15.7\%) & 1.5e-15 & 0.003 \\
Llama3.1:8B & Live write & GDM & 0.49/0.65/0.82 & 112 & 66 & 65 & 80 & 1.34 (16.4\%) & 7.6e-11 & 0.324 \\
Llama3.1:8B & Isolated, visible & GDM & 0.47/0.63/0.82 & 107 & 62 & 79 & 81 & 1.32 (16.2\%) & 2.1e-10 & 0.131 \\
\bottomrule\end{tabularx}
\end{table*}

